% Luc Destombe --- licence CC BY-NC-SA 4.0
% https://creativecommons.org/licenses/by-nc-sa/4.0/deed.fr
% Fichier autonome : compiler avec pdflatex, deux fois (signets, nombre de pages).
\def\lyceeversion{1.0}
\documentclass[a4paper,10pt]{article}
\makeatletter
\newif\iflycee@niveaux
\lycee@niveauxtrue
\newif\iflycee@cm
\newif\iflycee@palatino

\RequirePackage[utf8]{inputenc}
\RequirePackage[T1]{fontenc}
\RequirePackage[french]{babel}
\RequirePackage{etoolbox}

\newcommand{\ShorthandsOff}{\@ifundefined{shorthandoff}{}{\shorthandoff{;:!?}}}
\newcommand{\ShorthandsOn}{\@ifundefined{shorthandon}{}{\shorthandon{;:!?}}}
\ShorthandsOff
\AtBeginDocument{\ShorthandsOn}
\AtBeginEnvironment{tikzpicture}{\ShorthandsOff}

\RequirePackage{geometry}
\RequirePackage{fancyhdr}
\RequirePackage{lastpage}
\RequirePackage{multicol}
\RequirePackage{array}
\RequirePackage{enumitem}
\RequirePackage{xcolor}
\RequirePackage{needspace}

\RequirePackage{amsmath,amssymb}
\iflycee@cm\else\RequirePackage{fourier}\fi
\DeclareMathAlphabet{\mathcal}{OMS}{cmsy}{m}{n}
\renewcommand{\ttdefault}{lmtt}
\begingroup\catcode`\_=\active \catcode`\^=\active
\global\let_\sb \global\let^\sp
\endgroup
\AtBeginDocument{\catcode`\_=12 \mathcode`\_="8000
  \catcode`\^=12 \mathcode`\^="8000 }
\newlength{\retraitcalculs}
\setlength{\retraitcalculs}{2em}

\RequirePackage{tikz}
\usetikzlibrary{arrows.meta,calc,babel,shapes.misc}
\RequirePackage[most]{tcolorbox}
\tcbuselibrary{skins,breakable}

\definecolor{coulDef}{HTML}{1F4E79}
\definecolor{coulProp}{HTML}{7030A0}
\definecolor{coulRem}{HTML}{C55A11}
\definecolor{coulEx}{HTML}{2E7D32}
\definecolor{coulExo}{HTML}{B03A2E}
\definecolor{coulPart}{HTML}{1F4E79}
\definecolor{coulMeth}{HTML}{1B6E4F}
\definecolor{coulCourbe}{HTML}{0B5ED7}
\definecolor{coulCourbeB}{HTML}{C00000}
\definecolor{coulCourbeC}{HTML}{2E7D32}
\definecolor{coulGrille}{HTML}{8C8C8C}
\definecolor{coulRep}{HTML}{1B6E4F}
\definecolor{coulBar}{HTML}{2E7D32}

\newcommand{\R}{\mathbb{R}}
\newcommand{\Rs}{\mathbb{R}^{*}}
\newcommand{\N}{\mathbb{N}}
\newcommand{\Ns}{\mathbb{N}^{*}}
\newcommand{\Z}{\mathbb{Z}}
\newcommand{\Q}{\mathbb{Q}}
\renewcommand{\le}{\leqslant}
\renewcommand{\ge}{\geqslant}
\renewcommand{\approx}{\simeq}
\newcommand{\vect}[1]{\overrightarrow{#1}}
\newcommand{\norme}[1]{\left\lVert #1 \right\rVert}
\newcommand{\intervalle}[2]{\left[#1\,;#2\right]}
\RequirePackage{cancel}
\renewcommand{\CancelColor}{\color{coulExo}}
\newcommand{\fois}[1]{\times{\color{coulExo}#1}}
\newcommand{\barre}[1]{\cancel{#1}}
\newlength{\lycee@eqh}\newlength{\lycee@eqd}
\newsavebox{\lycee@eqa}\newsavebox{\lycee@eqb}\newsavebox{\lycee@eqc}
\newcommand{\lycee@eqcol}[1]{\hbox{$\renewcommand{\arraystretch}{1.35}%
  \begin{array}[t]{@{}>{\displaystyle}r@{}>{\displaystyle{}}l@{}}#1\end{array}$}}
\newcommand{\lycee@eqmesure}[1]{%
  \ifdim\ht#1>\lycee@eqh\setlength{\lycee@eqh}{\ht#1}\fi
  \ifdim\dp#1>\lycee@eqd\setlength{\lycee@eqd}{\dp#1}\fi}
\newcommand{\lycee@eqtrait}{\hspace{1.2em}\tikz[baseline=0pt]\draw[dash pattern=on 2pt off 2pt,
  line width=0.5pt] (0,\lycee@eqh) -- (0,-\lycee@eqd);\hspace{1.2em}}
\newcommand{\eqproduit}[2]{\par\smallskip\noindent C'est une équation produit :
  $A\times B=0 \iff A=0$ ou $B=0$.\par\nopagebreak\smallskip
  \begingroup\setlength{\lycee@eqh}{0pt}\setlength{\lycee@eqd}{0pt}%
  \setbox\lycee@eqa=\lycee@eqcol{#1}\setbox\lycee@eqb=\lycee@eqcol{#2}%
  \lycee@eqmesure\lycee@eqa\lycee@eqmesure\lycee@eqb%
  \noindent\hspace*{\retraitcalculs}\usebox\lycee@eqa\lycee@eqtrait\usebox\lycee@eqb\endgroup\par\smallskip}
\newcommand{\eqproduittrois}[3]{\par\smallskip\noindent C'est une équation produit :
  $A\times B\times C=0 \iff A=0$ ou $B=0$ ou $C=0$.\par\nopagebreak\smallskip
  \begingroup\setlength{\lycee@eqh}{0pt}\setlength{\lycee@eqd}{0pt}%
  \setbox\lycee@eqa=\lycee@eqcol{#1}\setbox\lycee@eqb=\lycee@eqcol{#2}\setbox\lycee@eqc=\lycee@eqcol{#3}%
  \lycee@eqmesure\lycee@eqa\lycee@eqmesure\lycee@eqb\lycee@eqmesure\lycee@eqc%
  \noindent\hspace*{\retraitcalculs}\usebox\lycee@eqa\lycee@eqtrait\usebox\lycee@eqb\lycee@eqtrait\usebox\lycee@eqc\endgroup
  \par\smallskip}

\tikzset{grille/.style={coulGrille,line width=0.4pt}}
\newcommand{\quadrillage}[4]{\draw[grille,step=1] (#1,#2) grid (#3,#4);}
\tikzset{
  croix/.style={cross out,draw,line width=0.6pt,inner sep=0pt,minimum size=4pt},
  croixdroite/.style={croix,rotate=45,line width=0.8pt},
}
\newcommand{\pt}[3]{\path (#1) node[croix]{} node[#2,font=\small,inner sep=2.5pt]{$#3$};}

\newcommand{\lycee@repere}[4]{%
  \draw[grille,step=1] (#1,#2) grid (#3,#4);
  \draw[-{Stealth[length=2mm]},thick] (#1,0)--({#3+0.4},0) node[below left=-1pt]{$x$};
  \draw[-{Stealth[length=2mm]},thick] (0,#2)--(0,{#4+0.4}) node[below left=-1pt]{$y$};
  \node[below left,font=\scriptsize,inner sep=1.5pt] at (0,0) {$0$};}
\newcommand{\lycee@gradx}[1]{\foreach \k/\l in {#1}{%
  \draw (\k,0.07)--(\k,-0.07) node[below,font=\scriptsize,inner sep=1.5pt]{$\l$};}}
\newcommand{\lycee@grady}[1]{\foreach \k/\l in {#1}{%
  \draw (0.07,\k)--(-0.07,\k) node[left,font=\scriptsize,inner sep=1.5pt]{$\l$};}}
\newcommand{\lycee@intff}[2]{\left[#1\,;#2\right]}
\newcommand{\lycee@intfo}[2]{\left[#1\,;#2\right[}
\newcommand{\lycee@intof}[2]{\left]#1\,;#2\right]}
\newcommand{\lycee@intoo}[2]{\left]#1\,;#2\right[}
\newcommand{\lycee@ens}[1]{\left\{#1\right\}}
\AtBeginDocument{%
  \providecommand{\repere}{\lycee@repere}%
  \providecommand{\gradx}{\lycee@gradx}%
  \providecommand{\grady}{\lycee@grady}%
  \providecommand{\Cf}{\mathcal{C}\sb{f}}%
  \providecommand{\Cg}{\mathcal{C}\sb{g}}%
  \providecommand{\intff}{\lycee@intff}%
  \providecommand{\intfo}{\lycee@intfo}%
  \providecommand{\intof}{\lycee@intof}%
  \providecommand{\intoo}{\lycee@intoo}%
  \providecommand{\ens}{\lycee@ens}}

\newlist{questions}{enumerate}{3}
\setlist[questions,1]{label=\arabic*\textdegree),leftmargin=*,itemsep=2pt,topsep=3pt}
\setlist[questions,2]{label=\alph*),leftmargin=*,itemsep=1pt,topsep=2pt}
\setlist[questions,3]{label=\roman*),leftmargin=*,itemsep=1pt,topsep=1pt}
\setlist[itemize]{label=\textbullet,leftmargin=*,itemsep=2pt,topsep=3pt}

\newcommand{\besoinplace}[1]{\par\penalty\@M\begingroup
  \dimen@=#1\relax\dimen@ii=\pagegoal\advance\dimen@ii-\pagetotal
  \ifdim\dimen@>\dimen@ii\ifdim\dimen@ii>\z@\vfil\fi\break\fi\endgroup}

\newcommand{\lycee@pied}{\thepage/\pageref{LastPage}}
\newcommand{\entete}[2]{%
  \pagestyle{fancy}\fancyhf{}%
  \fancyhead[L]{\footnotesize\color{coulPart}\textbf{#1}}%
  \fancyhead[R]{\footnotesize\color{coulPart}\textbf{#2}}%
  \fancyfoot[C]{\footnotesize\lycee@pied}%
  \renewcommand{\headrulewidth}{0.4pt}%
  \renewcommand{\headrule}{\hbox to\headwidth{%
    \color{coulPart!40}\leaders\hrule height \headrulewidth\hfill}}}

\RequirePackage{ccicons}
\newcommand{\lycee@licence}{{\scriptsize\color{gray}%
  \copyright\ Luc Destombe --- \ccbyncsa\ CC BY-NC-SA 4.0}}
\newif\iflycee@licence \lycee@licencetrue
\newcommand{\sanslicence}{\lycee@licencefalse}
\AtBeginDocument{\iflycee@licence\AddToHookNext{shipout/foreground}{%
  \put(0,0){\rlap{\kern\dimexpr1in+\hoffset+\oddsidemargin+\textwidth\relax
    \llap{\raisebox{-\dimexpr1in+\voffset+\topmargin+\headheight+\headsep
      +\textheight+\footskip\relax}{\lycee@licence}}}}}\fi}

\newcommand{\lycee@site}{faire-des-maths.fr}
\newcommand{\lycee@version}{{\scriptsize\color{gray}\lycee@site
  \ifdefined\lyceeversion\ --- Version~\lyceeversion\fi}}
\AtBeginDocument{\iflycee@licence\AddToHookNext{shipout/foreground}{%
  \put(0,0){\rlap{\kern\dimexpr1in+\hoffset+\oddsidemargin\relax
    \raisebox{-\dimexpr1in+\voffset+\topmargin+\headheight+\headsep
      +\textheight+\footskip\relax}{\lycee@version}}}}\fi}

\newcommand{\formulesserrees}{%
  \setlength{\abovedisplayskip}{3pt}\setlength{\belowdisplayskip}{3pt}%
  \setlength{\abovedisplayshortskip}{2pt}\setlength{\belowdisplayshortskip}{3pt}}

\setlength{\parindent}{0pt}

\RequirePackage[implicit=false,hidelinks,pdfencoding=auto,psdextra,bookmarksopen,
  bookmarksopenlevel=1,pdfcreator={},pdfproducer={}]{hyperref}
\RequirePackage{bookmark}
\hypersetup{pdfauthor={Luc Destombe},pdfkeywords={CC BY-NC-SA 4.0}}
\newcounter{lycee@signet}
\newcommand{\signet}[3][]{\stepcounter{lycee@signet}%
  \edef\lycee@dest{\if\relax\detokenize{#1}\relax
    signet.\arabic{lycee@signet}\else#1\fi}%
  \hypertarget{\lycee@dest}{}\bookmark[level=#2,dest=\lycee@dest]{#3}}
\pdfstringdefDisableCommands{%
  \def\R{R}\def\N{N}\def\Z{Z}\def\Q{Q}\def\vect#1{#1}\def\Cf{Cf}\def\Cg{Cg}\def\fois#1{×#1}%
  \def\barre#1{#1}\def\intervalle#1#2{[#1;#2]}\def\intff#1#2{[#1;#2]}%
  \def\textsuperscript#1{#1}\def\dfrac#1#2{#1/#2}\def\frac#1#2{#1/#2}%
  \def\vec#1{#1⃗}}%
\RequirePackage[official]{eurosym}

\geometry{top=0.8cm,bottom=1.3cm,left=1cm,right=1cm,footskip=0.6cm}

\pagestyle{fancy}
\fancyhf{}
\renewcommand{\headrulewidth}{0pt}
\fancyfoot[C]{\footnotesize Page \thepage{} / \pageref{LastPage}}

\newcommand{\titrefiche}[2]{%
  \begin{center}
    \Large\textbf{#1}\\[2pt]
    \large\textbf{#2}\\[2pt]
  \end{center}
  \vspace{0.10cm}}

\setlength{\columnseprule}{0.4pt}
\setlength{\columnsep}{15pt}
\renewcommand{\columnseprulecolor}{\color{black!45}}
\setlength{\multicolsep}{2pt}

\newcounter{partiefiche}
\renewcommand{\thepartiefiche}{\Roman{partiefiche}}
\newif\ifapresbandeau
\newcommand{\lycee@bandeau}[1]{%
  \par\penalty-600
  \vspace{0.08cm}%
  \noindent\colorbox{coulPart}{%
    \parbox{\dimexpr\linewidth-2\fboxsep\relax}{%
      \color{white}\bfseries\raggedright\strut#1\strut}}%
  \nopagebreak\par\nopagebreak
  \vspace{0.06cm}\nopagebreak
  \apresbandeautrue}
\newcommand{\partiefiche}[1]{%
  \refstepcounter{partiefiche}%
  \lycee@bandeau{\signet{1}{\thepartiefiche{} --- #1}\thepartiefiche{} --- #1}}
\newcommand{\partiefichesansnum}[1]{\lycee@bandeau{\signet{1}{#1}#1}}

\newcounter{exo}
\newlength{\espaceexo}\setlength{\espaceexo}{7pt}
\newlength{\espacetitre}\setlength{\espacetitre}{2pt}
\newcommand{\exo}[1][\relax]{%
  \par
  \ifapresbandeau\apresbandeaufalse\else\penalty-150\addvspace{\espaceexo}\fi
  \refstepcounter{exo}%
  \noindent\signet[exercice-\arabic{exo}]{2}{Exercice \arabic{exo}}%
  \textbf{\color{coulExo}\underline{Exercice \theexo}}%
  \ifx\relax#1\relax\else\textbf{ : #1}\fi
  \par\nopagebreak\vspace{\espacetitre}\nopagebreak
  \ignorespaces}

\setlist[enumerate]{nosep,leftmargin=*,itemsep=2pt,topsep=2pt}
\setlist[itemize]{nosep,leftmargin=*,topsep=2pt}
\newcommand{\choix}[3]{\par\hspace*{1em}%
  \makebox[0.3\linewidth][l]{a) #1}\makebox[0.3\linewidth][l]{b) #2}c) #3}
\newcommand{\deux}[2]{\par\hspace*{0.6em}%
  \makebox[0.48\linewidth][l]{#1}#2}
\newcommand{\trois}[3]{\par\hspace*{0.6em}%
  \makebox[0.32\linewidth][l]{#1}\makebox[0.32\linewidth][l]{#2}#3}

  \renewcommand{\exo}[1][\relax]{%
    \ifapresbandeau\apresbandeaufalse\else\par\penalty-150\fi
    \refstepcounter{exo}%
    \vspace{0.05cm}%
    \noindent\signet[exercice-\arabic{exo}]{2}{Exercice \arabic{exo}}%
    \textbf{\color{coulExo}\underline{Exercice \theexo}}%
    \ifx\relax#1\relax\else\textbf{ : #1}\fi%
    \par\nopagebreak
    \ignorespaces}
  \newcommand{\rappel}[1]{%
    \par\vspace{0.06cm}\nopagebreak
    \noindent\fcolorbox{black!35}{black!5}{%
      \parbox{\dimexpr\linewidth-2\fboxsep-2\fboxrule\relax}{%
        \small\textbf{Rappel.} #1}}%
    \par\vspace{0.06cm}\nopagebreak}
  \setlist[enumerate]{nosep,leftmargin=*}
  \setlist[itemize]{nosep,leftmargin=*}
  \setlength{\multicolsep}{3pt plus 1pt minus 1pt}
  \setlength{\premulticols}{10pt}
  \setlength{\postmulticols}{10pt}
  \newlist{expr}{enumerate}{1}
  \setlist[expr]{label={\Alph*\ $=$},nosep,leftmargin=*,itemsep=0pt}

\renewenvironment{center}{\par\addvspace{3pt}\centering}{\par\addvspace{3pt}}
\formulesserrees
\AtBeginDocument{\formulesserrees}
\clubpenalty=10000
\widowpenalty=10000
\displaywidowpenalty=10000
\brokenpenalty=10000
\setlength{\parskip}{0pt}

\RequirePackage{ragged2e}
\setlength{\RaggedRightRightskip}{0pt plus 1fil}
\AtBeginDocument{\RaggedRight\binoppenalty=10000 \relpenalty=10000 }
\g@addto@macro\@parboxrestore{\RaggedRight}
\makeatother
\usepackage{tkz-tab}
\newcommand{\qcm}[4]{\par\hspace*{0.6em}\makebox[0.48\linewidth][l]{a) #1}b) #2
  \par\hspace*{0.6em}\makebox[0.48\linewidth][l]{c) #3}d) #4}

\begin{document}

\titrefiche{1\textsuperscript{re} STMG --- Dérivation et étude de fonctions}{Fiche d'exercices du chapitre}

\begin{multicols*}{2}\raggedcolumns

\partiefiche{Fonctions affines et droites}

\exo[Coefficient directeur]
Pour chaque couple de points, calculer le coefficient directeur de la droite $(AB)$.
\begin{enumerate}[label=\alph*)]
  \item $A(1\,;2)$ et $B(3\,;8)$
  \item $A(-1\,;4)$ et $B(2\,;-2)$
  \item $A(0\,;-3)$ et $B(4\,;-1)$
  \item $A(-3\,;5)$ et $B(1\,;5)$
\end{enumerate}

\exo[Équation d'une droite]
Pour chaque couple de points, déterminer l'équation réduite de la droite $(AB)$.
\begin{enumerate}[label=\alph*)]
  \item $A(1\,;1)$ et $B(2\,;4)$
  \item $A(-1\,;5)$ et $B(1\,;1)$
  \item $A(-2\,;-3)$ et $B(2\,;5)$
  \item $A(2\,;0)$ et $B(6\,;2)$
\end{enumerate}

\exo[Une fonction affine par deux images]
\begin{enumerate}[label=\arabic*)]
  \item $f$ est affine, avec $f(2)=7$ et $f(-1)=-2$. Déterminer $f(x)$, puis calculer
  $f(10)$.
  \item $g$ est affine, avec $g(-3)=8$ et $g(1)=0$. Déterminer $g(x)$, puis calculer
  $g(5)$.
\end{enumerate}

\exo[Lire des équations]
Lire l'équation réduite de chacune des droites $(d_{1})$, $(d_{2})$ et $(d_{3})$.
\begin{center}
\begin{tikzpicture}[x=0.45cm,y=0.45cm]
  \repere{-3}{-3}{5}{4}
  \gradx{-2,-1,1,2,3,4} \grady{-2,-1,1,2,3}
  \begin{scope}
  \clip (-3,-3) rectangle (5,4);
  \draw[coulCourbe,line width=1.1pt] (-1,-3)--(2.5,4);
  \draw[coulCourbeB,line width=1.1pt] (-1,4)--(5,-2);
  \draw[coulCourbeC,line width=1.1pt] (-2,-3)--(5,0.5);
  \end{scope}
  \node[coulCourbe,font=\footnotesize] at (1.4,3.6) {$(d_{1})$};
  \node[coulCourbeB,font=\footnotesize] at (-0.6,2.7) {$(d_{2})$};
  \node[coulCourbeC,font=\footnotesize] at (4.2,0.8) {$(d_{3})$};
\end{tikzpicture}
\end{center}

\partiefiche{Taux de variation}

\exo[Calculer des taux de variation]
Soit $f(x)=x^{2}-2x+3$. Calculer le taux de variation de $f$ :
\begin{enumerate}[label=\alph*)]
  \item entre $0$ et $2$ ;
  \item entre $1$ et $4$ ;
  \item entre $-2$ et $1$.
\end{enumerate}

\exo[Les clients d'un site de vente en ligne]
Le nombre de clients d'un site de vente en ligne, en milliers, $t$ mois après son
ouverture, est donné par le tableau suivant.
\begin{center}
\renewcommand{\arraystretch}{1.4}
\begin{tabular}{|c|c|c|c|c|}
\hline
$t$ (mois) & $0$ & $2$ & $5$ & $10$\\
\hline
$N(t)$ (milliers) & $4$ & $10$ & $16$ & $21$\\
\hline
\end{tabular}
\end{center}
\begin{enumerate}[label=\alph*)]
  \item Calculer le taux de variation de $N$ entre $0$ et $2$, entre $2$ et $5$, puis
  entre $5$ et $10$.
  \item Interpréter chacun de ces taux par une phrase.
  \item Pendant quelle période le nombre de clients augmente-t-il le plus vite en
  moyenne ?
\end{enumerate}

\exo[Lire un taux de variation sur une courbe]
On a tracé la courbe $\Cf$ d'une fonction $f$.
\begin{center}
\begin{tikzpicture}[x=0.45cm,y=0.45cm]
  \repere{-4}{-2}{4}{6}
  \gradx{-3,-2,-1,1,2,3} \grady{-1,1,2,3,4,5}
  \draw[coulCourbe,line width=1.1pt]
    (-3,-1) .. controls (-2.333,1) and (-1.667,3) .. (-1,3)
            .. controls (-0.333,3) and (0.333,1) .. (1,1)
            .. controls (1.667,1) and (2.333,3) .. (3,5);
  \pt{-3,-1}{left}{}
  \pt{-1,3}{above}{}
  \pt{1,1}{below}{}
  \pt{3,5}{left}{}
  \node[coulCourbe,font=\footnotesize] at (2.9,3.4) {$\Cf$};
\end{tikzpicture}
\end{center}
\begin{enumerate}[label=\alph*)]
  \item Lire $f(-3)$, $f(-1)$, $f(1)$ et $f(3)$.
  \item Calculer le taux de variation de $f$ entre $-3$ et $-1$, puis entre $-1$ et $1$.
  \item Tracer la sécante qui passe par les points de $\Cf$ d'abscisses $-3$ et $3$, puis
  lire son coefficient directeur. Que représente ce nombre pour $f$ ?
\end{enumerate}

\exo[Un coût de production]
Un atelier fabrique $q$ lampes par mois, avec $0\le q\le 50$. Le coût de production,
en euros, est $C(q)=q^{2}+20q+100$.
\begin{enumerate}[label=\alph*)]
  \item Calculer $C(10)$, $C(20)$ et $C(30)$.
  \item Calculer le taux de variation de $C$ entre $10$ et $20$, puis l'interpréter.
  \item Même question entre $20$ et $30$.
  \item Le coût augmente-t-il de plus en plus vite ?
  \item Déterminer l'équation réduite de la droite qui passe par les points de
  coordonnées $(10\,;400)$ et $(20\,;900)$.
\end{enumerate}

\partiefiche{Nombre dérivé}

\exo[Nombre dérivé en $1$]
Soit $f(x)=x^{2}+2x-1$.
\begin{enumerate}[label=\alph*)]
  \item Calculer $f(1)$, puis $f(1+h)$.
  \item Simplifier le taux de variation $T(h)=\dfrac{f(1+h)-f(1)}{h}$.
  \item En déduire $f'(1)$.
\end{enumerate}

\exo[Nombre dérivé en $2$]
Soit $g(x)=2x^{2}-3x$. Calculer $g(2)$ et $g(2+h)$, simplifier
le taux de variation $T(h)=\dfrac{g(2+h)-g(2)}{h}$, puis en déduire $g'(2)$.

\exo[Avec un coefficient négatif]
Soit $k(x)=-x^{2}+5x$. Calculer $k(1)$ et $k(1+h)$, simplifier
le taux de variation $T(h)=\dfrac{k(1+h)-k(1)}{h}$, puis en déduire $k'(1)$.

\exo[Vers le nombre dérivé]
Soit $f(x)=x^{2}$. On donne $4{,}1^{2}=16{,}81$ et $4{,}01^{2}=16{,}0801$.
\begin{enumerate}[label=\alph*)]
  \item Compléter le tableau, où $T(h)=\dfrac{f(4+h)-f(4)}{h}$.
  \begin{center}
  \renewcommand{\arraystretch}{1.4}
  \begin{tabular}{|c|c|c|c|}
  \hline
  $h$ & $1$ & $0{,}1$ & $0{,}01$\\
  \hline
  $T(h)$ & & & \\
  \hline
  \end{tabular}
  \end{center}
  \item Quelle valeur de $f'(4)$ ce tableau suggère-t-il ?
\end{enumerate}

\partiefiche{Tangente à une courbe}

\exo[Lire des nombres dérivés]
\begin{minipage}[c]{0.42\linewidth}\raggedright
On a tracé $\Cf$ et ses tangentes aux points $A$, $B$, $C$ et $E$.

\smallskip
Lire $f'(-3)$, $f'(-1)$, $f'(1)$ et $f'(4)$.
\end{minipage}\hfill
\begin{minipage}[c]{0.56\linewidth}\centering
\begin{tikzpicture}[x=0.45cm,y=0.45cm]
  \repere{-4}{-2}{5}{4}
  \gradx{-3,-2,-1,1,2,3,4} \grady{-1,1,2,3}
  \draw[coulCourbe,line width=1.1pt]
    (-3,2) .. controls (-2.333,0.667) and (-1.667,-1) .. (-1,-1)
           .. controls (-0.333,-1) and (0.333,-0.333) .. (1,1)
           .. controls (1.667,2.333) and (2.333,3) .. (3,3)
           .. controls (3.333,3) and (3.667,2.667) .. (4,2);
  \draw[coulCourbeB,thick] (-3.8,3.6)--(-2.2,0.4);
  \draw[coulCourbeB,thick] (-2,-1)--(0,-1);
  \draw[coulCourbeB,thick] (0.2,-0.6)--(1.8,2.6);
  \draw[coulCourbeB,thick] (3.4,3.2)--(4.6,0.8);
  \pt{-3,2}{left}{A}
  \pt{-1,-1}{below}{B}
  \pt{1,1}{right}{C}
  \pt{4,2}{right}{E}
  \node[coulCourbe,font=\footnotesize] at (2.2,3.5) {$\Cf$};
\end{tikzpicture}
\end{minipage}

\exo[Tracer des tangentes]
Dans un repère, placer chaque point de $\Cf$, puis tracer la tangente à $\Cf$ en ce
point.
\begin{enumerate}[label=\alph*)]
  \item $A(0\,;1)$ et $f'(0)=2$
  \item $B(1\,;3)$ et $f'(1)=-1$
  \item $C(-2\,;-1)$ et $f'(-2)=\frac{1}{2}$
  \item $D(3\,;-2)$ et $f'(3)=0$
\end{enumerate}

\exo[Une tangente qui passe par deux points]
Pour chacune des tangentes $(T_{1})$, $(T_{-1})$ et $(T_{3})$ à $\Cf$, donner $f(a)$ et
calculer $f'(a)$, où $a$ est l'abscisse du point de contact.
\begin{enumerate}[label=\alph*)]
  \item $(T_{1})$ touche $\Cf$ au point $A(1\,;2)$ et passe par $B(3\,;6)$.
  \item $(T_{-1})$ touche $\Cf$ au point $C(-1\,;4)$ et passe par $D(2\,;-2)$.
  \item $(T_{3})$ touche $\Cf$ au point $E(3\,;5)$ et est parallèle à l'axe des
  abscisses.
\end{enumerate}

\exo[Tangentes qui montent, tangentes qui descendent]
On a tracé la courbe $\Cf$ d'une fonction $f$.
\begin{center}
\begin{tikzpicture}[x=0.45cm,y=0.45cm]
  \repere{-5}{-2}{4}{4}
  \gradx{-4,-3,-2,-1,1,2,3} \grady{-1,1,2,3}
  \draw[coulCourbe,line width=1.1pt]
    (-4,-1) .. controls (-3.333,1) and (-2.667,3) .. (-2,3)
            .. controls (-1,3) and (0,-1) .. (1,-1)
            .. controls (1.667,-1) and (2.333,1) .. (3,3);
  \node[coulCourbe,font=\footnotesize] at (1.6,2.4) {$\Cf$};
\end{tikzpicture}
\end{center}
\begin{enumerate}[label=\alph*)]
  \item En quels points de $\Cf$ la tangente est-elle horizontale ? Que valent les
  nombres dérivés en ces points ?
  \item Sans tracer de tangente, donner le signe de $f'(-3)$, de $f'(0)$ et de $f'(2)$.
  \item Écrire une phrase qui relie le nombre $f'(-2)$ et la tangente à $\Cf$ au point
  d'abscisse $-2$.
\end{enumerate}

\exo[Retrouver l'allure d'une courbe]
Dans un repère, tracer une courbe qui passe par $A(-3\,;-2)$, $B(-1\,;2)$,
$C(1\,;0)$ et $D(3\,;-2)$, avec $f'(-3)=3$, $f'(-1)=0$, $f'(1)=-2$ et $f'(3)=0$.
Commencer par tracer les quatre tangentes.

\partiefiche{Équation de la tangente}

\exo[Avec un nombre dérivé admis]
Soit $f(x)=x^{2}-3x+2$. On admet que $f'(4)=5$ et $f'(1)=-1$.
\begin{enumerate}[label=\alph*)]
  \item Calculer $f(4)$, puis déterminer l'équation réduite de la tangente à $\Cf$
  au point d'abscisse $4$.
  \item Même question au point d'abscisse $1$.
\end{enumerate}

\exo[Quatre tangentes]
Pour chaque cas, déterminer l'équation réduite de la tangente à $\Cf$ au point
d'abscisse $a$.
\begin{enumerate}[label=\alph*)]
  \item $a=2$ ; $f(2)=3$ et $f'(2)=1$
  \item $a=-1$ ; $f(-1)=4$ et $f'(-1)=-2$
  \item $a=3$ ; $f(3)=-1$ et $f'(3)=0$
  \item $a=0$ ; $f(0)=5$ et $f'(0)=-3$
\end{enumerate}

\exo[Lire, puis calculer]
On a tracé la courbe $\Cf$ et sa tangente $(T_{-1})$ au point $A$ d'abscisse $-1$.
\begin{center}
\begin{tikzpicture}[x=0.5cm,y=0.5cm]
  \repere{-3}{-3}{3}{4}
  \gradx{-2,-1,1,2} \grady{-2,-1,1,2,3}
  \begin{scope}
  \clip (-3,-3) rectangle (3,4);
  \draw[coulCourbe,line width=1.1pt,domain=-2.4:2.4,samples=50] plot (\x,{-\x*\x+2});
  \draw[coulCourbeB,thick] (-3,-3)--(0.5,4);
  \end{scope}
  \pt{-1,1}{left}{A}
  \node[coulCourbe,font=\footnotesize] at (1.8,0.9) {$\Cf$};
  \node[coulCourbeB,font=\footnotesize] at (1.2,3.5) {$(T_{-1})$};
\end{tikzpicture}
\end{center}
\begin{enumerate}[label=\alph*)]
  \item Lire $f(-1)$ et $f'(-1)$.
  \item En déduire l'équation réduite de $(T_{-1})$.
  \item Lire le point où $(T_{-1})$ coupe l'axe des ordonnées, et contrôler la réponse
  b).
\end{enumerate}

\exo[Des équations à contrôler]
Pour chaque affirmation, dire si elle est vraie ou fausse, en justifiant.
\begin{enumerate}[label=\alph*)]
  \item Soit $f(x)=x^{2}-2x+4$ ; on admet que $f'(3)=4$. \emph{La tangente à $\Cf$ au
  point d'abscisse $3$ a pour équation $y=4x-5$.}
  \item Soit $g(x)=-x^{2}+5$ ; on admet que $g'(1)=-2$. \emph{La tangente à $\Cg$ au
  point d'abscisse $1$ a pour équation $y=-2x+4$.}
  \item On sait que $h(2)=5$ et $h'(2)=0$. \emph{La tangente à la courbe de $h$ au point
  d'abscisse $2$ a pour équation $y=5$.}
\end{enumerate}

\exo[Du taux de variation à la tangente]
Soit $f(x)=-x^{2}+3x+1$.
\begin{enumerate}[label=\alph*)]
  \item Calculer $f(2)$, puis $f(2+h)$.
  \item Simplifier le taux de variation $T(h)=\dfrac{f(2+h)-f(2)}{h}$ et en déduire
  $f'(2)$.
  \item Déterminer l'équation réduite de la tangente à $\Cf$ au point d'abscisse $2$.
\end{enumerate}

\partiefiche{Fonction dérivée d'une fonction polynôme}

\exo[Dériver]
Calculer la dérivée de chaque fonction.
\begin{multicols}{2}
\begin{enumerate}[label=\alph*)]
  \item $f(x)=7x-5$
  \item $g(x)=1-4x$
  \item $h(x)=x^{2}+6x-2$
  \item $k(x)=3x^{2}-x+8$
  \item $m(x)=-5x^{2}+10x$
  \item $p(x)=x^{3}+4x^{2}-x$
  \item $q(x)=2x^{3}-3x^{2}+1$
  \item $r(x)=-2x^{3}+5$
\end{enumerate}
\end{multicols}

\exo[Entraînement]
Calculer la dérivée de chaque fonction.
\begin{multicols}{2}
\begin{enumerate}[label=\alph*)]
  \item $f(x)=6x-2$
  \item $g(x)=-x^{2}+4x-3$
  \item $h(x)=5x^{3}-2x^{2}$
  \item $k(x)=x^{3}+x^{2}+x+1$
  \item $m(x)=-3x^{3}+6x$
  \item $p(x)=2x^{3}-5x^{2}+3$
  \item $q(x)=\frac{1}{2}x^{2}+3x$
  \item $r(x)=\frac{1}{3}x^{3}-x$
\end{enumerate}
\end{multicols}

\exo[Développer d'abord]
Développer chaque expression, puis calculer sa dérivée.
\begin{multicols}{2}
\begin{enumerate}[label=\alph*)]
  \item $A(x)=(x-1)(x+4)$
  \item $B(x)=(2x+1)(x-3)$
  \item $C(x)=(x+2)^{2}$
  \item $D(x)=x(x^{2}-5)$
\end{enumerate}
\end{multicols}

\exo[Dérivée et tangentes]
Soit $f(x)=x^{3}-2x+1$.
\begin{enumerate}[label=\alph*)]
  \item Calculer $f'(x)$.
  \item Calculer $f(1)$ et $f'(1)$, puis déterminer l'équation de la tangente au point
  d'abscisse $1$.
  \item Même question au point d'abscisse $0$.
\end{enumerate}

\partiefiche{Dérivée et sens de variation}

\exo[Du signe de $f'$ aux variations]
Voici le signe de $f'(x)$ sur $\intff{-4}{5}$. On sait que $f(-4)=2$, $f(-1)=-3$,
$f(3)=4$ et $f(5)=1$.
\begin{center}
\begin{tikzpicture}[scale=0.8]
  \tkzTabInit[lgt=1.6,espcl=1.8]{$x$/1,$f'(x)$/1}{$-4$,$-1$,$3$,$5$}
  \tkzTabLine{,-,z,+,z,-,}
\end{tikzpicture}
\end{center}
Dresser le tableau de variations de $f$.

\exo[Un trinôme]
Soit $f(x)=x^{2}-6x+5$ sur $\intff{0}{6}$.
\begin{enumerate}[label=\alph*)]
  \item Calculer $f'(x)$ et étudier son signe.
  \item Dresser le tableau de variations de $f$.
\end{enumerate}

\exo[Un polynôme de degré 3]
Soit $f(x)=x^{3}-3x^{2}+1$ sur $\intff{-1}{3}$.
\begin{enumerate}[label=\alph*)]
  \item Calculer $f'(x)$.
  \item Vérifier que $f'(x)=3x(x-2)$.
  \item Étudier le signe de $f'(x)$.
  \item Dresser le tableau de variations de $f$.
\end{enumerate}

\exo[Un coefficient négatif]
Soit $g(x)=-2x^{3}+3x^{2}+12x-5$ sur $\intff{-2}{3}$.
\begin{enumerate}[label=\alph*)]
  \item Calculer $g'(x)$.
  \item Vérifier que $g'(x)=-6(x+1)(x-2)$.
  \item Étudier le signe de $g'(x)$.
  \item Dresser le tableau de variations de $g$.
\end{enumerate}

\exo[Étude complète et tracé]
Soit $f(x)=x^{3}-3x$ sur $\intff{-2}{2}$.
\begin{enumerate}[label=\alph*)]
  \item Calculer $f'(x)$, puis vérifier que $f'(x)=3(x+1)(x-1)$.
  \item Étudier le signe de $f'(x)$, puis dresser le tableau de variations de $f$.
  \item Déterminer l'équation de la tangente $(T_{0})$ au point d'abscisse $0$.
  \item Compléter un tableau de valeurs de $f$ pour $x=-2$, $-1$, $0$, $1$ et $2$.
  \item Dans un repère, tracer $(T_{0})$, puis $\Cf$.
\end{enumerate}

\partiefiche{Problèmes concrets : maximum et minimum}

\exo[Un coût minimal]
Une usine fabrique $q$ dizaines de pièces par jour, avec $0\le q\le 25$. Le coût de
fabrication, en euros, est $C(q)=q^{2}-30q+300$.
\begin{enumerate}[label=\alph*)]
  \item Calculer $C'(q)$.
  \item Dresser le tableau de variations de $C$ sur $\intff{0}{25}$.
  \item Combien de pièces faut-il fabriquer pour que le coût soit minimal ? Que
  vaut-il ?
\end{enumerate}

\exo[La boîte la plus grande]
Une entreprise d'emballage fabrique des boîtes sans couvercle : dans un carton carré
de $12$~cm de côté, on découpe aux quatre coins un carré de côté $x$ (en cm), puis on
relève les bords. Le volume de la boîte est $V(x)=x(12-2x)^{2}$, pour $0\le x\le 6$.
\begin{enumerate}[label=\alph*)]
  \item Montrer que $V(x)=4x^{3}-48x^{2}+144x$.
  \item Calculer $V'(x)$, puis vérifier que $V'(x)=12(x-2)(x-6)$.
  \item Étudier le signe de $V'(x)$ sur $\intff{0}{6}$.
  \item Dresser le tableau de variations de $V$.
  \item Pour quelle valeur de $x$ le volume est-il maximal ? Que vaut-il ?
\end{enumerate}

\exo[Les inscriptions à une course solidaire]
Une association ouvre les inscriptions à sa course solidaire. Au bout de $t$ jours,
$N(t)=-0{,}5t^{2}+20t$ personnes sont inscrites, pour $0\le t\le 20$. Le nombre
$N'(t)$ est le rythme des inscriptions, en inscriptions par jour.
\begin{enumerate}[label=\alph*)]
  \item Calculer $N'(t)$.
  \item Quel est le rythme des inscriptions au début ($t=0$) ? Au bout de $10$ jours ?
  \item Au bout de combien de jours les inscriptions s'arrêtent-elles ?
  \item Combien de personnes sont alors inscrites ?
\end{enumerate}

\partiefichesansnum{Exercices type épreuve anticipée}

\exo[QCM]
Pour chaque question, une seule réponse est exacte.
\begin{enumerate}[label=\arabic*.]
  \item La dérivée de $f(x)=x^{3}-2x^{2}+5$ est :
        \qcm{$3x^{2}-4x+5$}{$3x^{2}-4x$}{$x^{2}-4x$}{$3x^{2}-2x$}
  \item Soit $g(x)=x^{2}-6x+3$. Le nombre dérivé $g'(1)$ vaut :
        \qcm{$-4$}{$-2$}{$4$}{$-5$}
  \item La tangente à $\Cf$ au point d'abscisse $2$ a pour équation $y=-3x+4$. Alors :
        \qcm{$f(2)=4$}{$f'(2)=4$}{$f'(2)=-3$}{$f(2)=-3$}
  \item Sur $\intff{0}{5}$, $f'(x)$ est négatif sur $\intff{0}{3}$ et positif sur
        $\intff{3}{5}$. La fonction $f$ :
        \qcm{croît sur $\intff{0}{5}$}{a un maximum en $3$}{a un minimum en $3$}{décroît
        sur $\intff{0}{5}$}
  \item Si $h'(x)=-2x+6$, la fonction $h$ est croissante sur :
        \qcm{$\left]-\infty\,;3\right]$}{$\left[3\,;+\infty\right[$}{$\left]-\infty\,;6\right]$}{$\left[-2\,;+\infty\right[$}
  \item Une fonction dont la dérivée est $3x^{2}+2$ est :
        \qcm{$x^{3}+2x$}{$x^{3}+2$}{$3x^{3}+2x$}{$6x$}
\end{enumerate}

\exo[Vrai ou faux ?]
Pour chacune des trois affirmations, dire si elle est vraie ou fausse, en justifiant.
\begin{enumerate}[label=\arabic*.]
  \item Soit $f(x)=4x^{2}-3x+1$. \emph{La tangente à $\Cf$ au point d'abscisse $1$ a
        pour équation $y=5x-3$.}
  \item Soit $g(x)=x^{3}-12x$. \emph{La fonction $g$ est croissante sur
        $\intff{0}{4}$.}
  \item Soit $h(x)=-x^{2}+8x$ sur $\intff{0}{8}$. \emph{Le maximum de $h$ sur
        $\intff{0}{8}$ vaut $16$.}
\end{enumerate}

\exo[Une courbe et sa tangente]
La courbe $\Cf$ ci-dessous représente la fonction $f$ définie sur $\intff{-2}{6}$ par
$f(x)=-0{,}5x^{2}+2x+4$. La droite $(T_{4})$ est la tangente à $\Cf$ au point
d'abscisse $4$ ; elle passe par les points de coordonnées $(4\,;4)$ et $(6\,;0)$.
\begin{center}
\begin{tikzpicture}[x=0.42cm,y=0.42cm]
  \repere{-3}{-3}{7}{7}
  \gradx{-2,-1,1,2,3,4,5,6} \grady{-2,-1,1,2,3,4,5,6}
  \begin{scope}
  \clip (-3,-3) rectangle (7,7);
  \draw[coulCourbe,line width=1.1pt,domain=-2:6,samples=50] plot (\x,{-0.5*\x*\x+2*\x+4});
  \draw[coulCourbeB,thick] (2.5,7)--(6.5,-1);
  \end{scope}
  \pt{4,4}{above right}{}
  \pt{6,0}{above right}{}
  \node[coulCourbe,font=\footnotesize] at (-1.2,2.4) {$\Cf$};
  \node[coulCourbeB,font=\footnotesize] at (6.3,1.6) {$(T_{4})$};
\end{tikzpicture}
\end{center}
\begin{enumerate}[label=\arabic*.]
  \item Par lecture graphique, donner $f(4)$ et $f'(4)$.
  \item Calculer $f'(x)$.
  \item Retrouver $f'(4)$ par le calcul.
  \item Déterminer l'équation réduite de $(T_{4})$.
  \item Étudier le signe de $f'(x)$ sur $\intff{-2}{6}$, puis dresser le tableau de
        variations de $f$.
  \item \emph{Affirmation : pour tout nombre réel $x$ de $\intff{-2}{6}$ on a
        $f(x)\le 6$.} Est-elle vraie ou fausse ?
\end{enumerate}

\exo[Le bénéfice d'un fabricant de vélos]
Une entreprise fabrique et vend chaque année $x$ centaines de vélos électriques, avec
$0\le x\le 10$. Son bénéfice, en milliers d'euros, est
$B(x)=-x^{3}+9x^{2}+48x-100$.
\begin{enumerate}[label=\arabic*.]
  \item Calculer $B(0)$ et interpréter ce résultat.
  \item Calculer $B'(x)$.
  \item Montrer que $B'(x)=(-3x+24)(x+2)$.
  \item Étudier le signe de $B'(x)$ sur $\intff{0}{10}$, puis dresser le tableau de
        variations de $B$. On donne $B(8)=348$ et $B(10)=280$.
  \item Combien de vélos l'entreprise doit-elle vendre pour que son bénéfice soit
        maximal ? Quel est ce bénéfice ?
\end{enumerate}

\end{multicols*}

\end{document}
